\section{本章小结}
\label{sec:mvd-summary}

本章从试验机制和取值范围出发，写出了五种常见的多元分布，并追问各自的依赖从哪里来。四种构造方式贯穿全章：独立试验的计数，标准高斯向量的仿射像，高斯向量乘以共同的随机尺度，独立Gamma变量除以总量。

计数模型来自独立试验或不放回抽样。多项分布由逐次独立、单次互斥得到协方差$n(\operatorname{diag}\bm{p}-\bm{p}\bm{p}^{\top})$；多元超几何分布的均值相同，协方差多乘有限总体因子$(M-n)/(M-1)$，抽样比例小时可以用多项分布近似。总数固定使协方差矩阵奇异，各类计数两两负相关。

多元高斯定义为$\bm{\mu}+\bm{A}\bm{Z}$的分布，它只由$\bm{\mu}$与$\bm{\Sigma}=\bm{A}\bm{A}^{\top}$决定，在仿射变换下封闭，边缘只需取子块。$\bm{\Sigma}$正定时有满维密度，等密度面是马氏距离的椭球；$\bm{\Sigma}$奇异时分布落在低维平面上。分块后，条件均值是观测的仿射函数，条件协方差是与观测无关的Schur补；残差与观测独立，这一正交分解也给出了高斯样本平均与残差平方和的独立性。零交叉协方差推出独立，只在联合高斯时成立。

多元t与多元Laplace把高斯的固定尺度换成共同的随机尺度：t用$\sqrt{\nu/U}$，尾部按负幂衰减，只有低于$\nu$阶的矩，$\bm{\Sigma}$是尺度矩阵；Laplace用$\sqrt W$，边缘是一元Laplace，各阶矩有限，$\bm{\Sigma}$就是协方差。两者都是椭圆分布，尺度矩阵为对角时分量不相关却不独立，大值倾向于同时出现；改用各分量独立的尺度，得到边缘与协方差相同而联合规律不同的乘积模型。

Dirichlet分布描述单纯形上的随机比例，密度相对于$d-1$个自由坐标给出，参数小于一时在边界附近无界却没有点质量。它可以由同速率独立Gamma变量归一化得到，边缘是Beta分布，均值为$\bm{\alpha}/\alpha_0$，协方差与$\alpha_0+1$成反比，并因总和为$1$而奇异。

本章多次用到“已知输入的分布，求输出的分布”：仿射像的密度、尺度混合的密度、Gamma归一化的联合密度。第~\ref{chap:random-variable-transformations}章把这些具体计算整理成一般的变换方法，补全了本章留给它的换元证明。
